\section{Discussion}
\label{sec:discussion}

\subsection{What the Empirical Study Changes}

The empirical validation sharpens three claims. First, the cost of
formal correctness is smaller than the analytical model conservatively
assumed: validation and selection measure in microseconds
(Section~\ref{subsec:empirical}, Finding~1), so the earlier
``below 3\% planning overhead'' bound is dominated entirely by optional
calibration, not by the algorithm. Second, the risk \gate{} guards
against is demonstrably real: a production compiler release silently
miscompiled a standard vision model in a way that per-operator testing
did not catch (Finding~3). Third, correctness-first selection has
observable consequences (Finding~2), which motivates the normalization
discussion below.

\subsection{Cost Normalization Sensitivity}

Max-normalization makes each metric dimension's \emph{largest} candidate
value the unit. When a dimension's spread is tiny---e.g.\ $m_5$
differing by $10^{-6}$ between exact and BatchNorm-folded plans---the
normalized difference is still $1.0$, and the dimension can dominate
selection despite being practically insignificant. This is
conservative (exactness-preserving) but can forgo real performance.
A principled alternative is \emph{threshold-relative} normalization for
bounded metrics: normalize $m_5$ by the constraint tolerance
$\epsilon$ rather than the candidate maximum, so deviations far below
tolerance contribute proportionally little. We leave a systematic
study of normalization schemes to future work; the linear cost model
itself (Lemma~\ref{lem:monotonicity}) is agnostic to this choice.

\subsection{Limitations}

\begin{itemize}[leftmargin=1.6em]
  \item \textbf{Candidate-set relativity.} Optimality
    (Theorem~\ref{thm:optimality}) is relative to
    $\mathcal{P}_{\text{cand}}$; heuristic generation may omit the
    globally optimal strategy.
  \item \textbf{Analytical error model.} Interval-arithmetic bounds may
    overestimate error and reject acceptable plans; the empirical
    implementations sidestep this via calibration, which trades
    planning time for tightness.
  \item \textbf{Single-device execution.} No distributed or multi-GPU
    execution yet.
  \item \textbf{Backend coverage.} CPU, CUDA, and Vulkan today; TPU
    access in this paper is via XLA/JAX rather than a native backend.
  \item \textbf{Energy metrology.} $m_4$ requires platform power
    counters; on the empirical CPU platform (Windows, no RAPL access)
    it degrades to an analytical estimate, and cloud TPU guests expose
    no power telemetry at all. The metric is well-defined but its
    \emph{measurement} is platform-gated.
  \item \textbf{Reference-implementation performance.} The
    dependency-free Rust reference is $9\times$ slower than ONNX
    Runtime on CPU; production deployments would pair \gate{} with
    optimized kernels (glproc/glcuda), for which the analytical study
    models the expected regime.
\end{itemize}

\subsection{Future Work}

(1) ML-guided plan generation to improve candidate quality; (2)
distributed execution across devices and nodes; (3) runtime constraint
adaptation under changing system load; (4) additional backends (native
TPU, WebGPU, SYCL); (5) drop-in runtime integration for PyTorch and
TensorFlow; (6) formal verification of the constraint implementations
themselves; and (7) threshold-relative cost normalization
(Section~\ref{sec:cost}).
